% =====================================================================
%  RESEARCH PAPER  --  economics article skeleton
%  Build:  latexmk -cd paper/paper-template.tex   (or the VSCode green arrow)
%
%  Bibliography: biblatex+biber by default. For a journal that requires
%  bibtex/natbib, change the line below to \usepackage[natbib]{meconbib}
%  and (optionally) \setbibstyle{aer}.  Nothing else in the file changes.
% =====================================================================
\documentclass[12pt]{article}

\usepackage{mathprelude}
\usepackage{setspace}
\usepackage{meconbib}       % <-- add [natbib] here for older journals
\bibresource{references}

% --- front-matter helpers ---
\newcommand{\keywords}[1]{\par\smallskip\noindent\textbf{Keywords:} #1}
\newcommand{\jelcodes}[1]{\par\noindent\textbf{JEL Classification:} #1}

\title{\textbf{An Informative Title About an Economic Question}%
  \thanks{I thank seminar participants for helpful comments. All errors are my own.}}
\author{%
  Gavin Qu\thanks{Department of Economics; \texttt{gavinqu1@gmail.com}.}%
}
\date{\today \\[6pt] \small\emph{Preliminary --- please do not circulate.}}

\begin{document}
\maketitle

\begin{abstract}\noindent
This paper studies [question]. Using [method], I show [main result]. The findings
imply [contribution]. \keywords{matching, mechanism design, information}
\jelcodes{C78, D47, D82}
\end{abstract}

\bigskip
% \doublespacing   % <- uncomment for a submission draft

\section{Introduction}\label{sec:intro}
Motivation, contribution, and roadmap. This paper relates to the classic analysis
of adverse selection by \textcite{akerlof1970}; for textbook foundations see
\textcite{mwg1995}.

\section{Model}\label{sec:model}
\begin{assumption}\label{as:reg}
Preferences are complete, transitive, continuous, and locally non-satiated.
\end{assumption}

\begin{proposition}\label{prop:exist}
Under \Cref{as:reg}, a competitive equilibrium exists.
\end{proposition}
\begin{proof}
See \Cref{app:proofs}.
\end{proof}

\section{Results}\label{sec:results}
Main analysis and \Cref{prop:exist}'s implications.

\section{Conclusion}\label{sec:conc}
Summary and directions for future work.

\printrefs

\appendix
\section{Proofs}\label{app:proofs}
\begin{proof}[Proof of \Cref{prop:exist}]
Apply a fixed-point argument to the excess-demand correspondence. \qedhere
\end{proof}

\end{document}
